% CVPR 2027 workshop submission: what visual imagination buys a program-writing agent.
% Kit: cvpr.sty and ieeenat_fullname.bst copied from ../cvpr27_faces (2027 kit).
% Build: `make` (draft) or `make submit` (fails on any leftover \TODO, \UNMEASURED or \com).
%
% Writing convention (Foerster, "How to write an ML paper"): every paragraph opens
% with a one-line `%TL;DR` comment that states its single idea. Read only the
% TL;DR lines to review the flow of the paper (`make tldr`).
%
% Ownership: teaser, 0_abstract, 1_intro, 2_related, 3_method, preamble and
% references.bib are the front half; 4_setup, 5_results, 6_conclusion and
% references_new.bib are written separately.
% Every number carries a LaTeX comment with its results.md entry (T number) or card.

\documentclass[10pt,twocolumn,letterpaper]{article}

%%%%%%%%% PAPER TYPE
% \usepackage{cvpr}              % CAMERA-READY
\usepackage[review]{cvpr}      % REVIEW version
% \usepackage[pagenumbers]{cvpr} % arXiv

%% Preamble (CVPR 2027 kit). Loaded before hyperref.
%% cvpr.sty already provides times, xcolor, graphicx, amsmath, amssymb, booktabs,
%% natbib[numbers], caption, subcaption, url, enumitem and cleveref.

\usepackage{multirow}
\usepackage{csquotes}   % \enquote{...}
\usepackage{algorithm}
\usepackage{algpseudocode}

\emergencystretch=1.5em

% ---------------------------------------------------------------------------
% Draft switch. `make submit` defines \submission; then every leftover \TODO,
% \UNMEASURED or \com is a hard error.
% ---------------------------------------------------------------------------
\newif\ifdraft\drafttrue
\ifdefined\submission\draftfalse\fi
\ifdraft
  \newcommand{\TODO}[1]{\textcolor{red}{\textbf{[TODO: #1]}}}
  \newcommand{\UNMEASURED}[1]{\textcolor{red}{\textbf{[UNMEASURED: #1]}}}
  \newcommand{\com}[1]{\textcolor{blue}{Comment: #1}}
\else
  \newcommand{\TODO}[1]{\PackageError{paper}{Unresolved TODO: #1}{}}
  \newcommand{\UNMEASURED}[1]{\PackageError{paper}{Unmeasured claim: #1}{}}
  \newcommand{\com}[1]{\PackageError{paper}{Unaddressed comment: #1}{}}
\fi
% Short form of \UNMEASURED for narrow table cells: a red "?" in draft, an error at submit.
\ifdraft
  \newcommand{\tbd}{\textcolor{red}{\textbf{?}}}
\else
  \newcommand{\tbd}{\PackageError{paper}{Unmeasured table cell}{}}
\fi

% ---------------------------------------------------------------------------
% Work-specific terminology. One name per concept, never a synonym.
%   build program   the text the builder LLM writes (commands of the build language)
%   build           the voxel structure a program produces on the 32^3 canvas
%   dream           the target image the image model draws from the prompt
%   caption         the structured description of a dream written by the captioner
%   spec            a structured description written from the prompt alone (no image)
%   builder         the LLM role that writes build programs
%   captioner       the LLM role that writes captions (never sees the prompt)
%   render          an image of a build from the textured renderer
%   match turn      one compare-and-repair turn of the dream loop
%   critique turn   one render-and-repair turn of the no-image loop
%   arm             one recipe evaluated on the whole prompt set
% Arms are small caps, metrics roman.
% ---------------------------------------------------------------------------
\newcommand{\classicm}{\mbox{\textsc{Classic-Matched}}}   % no image, 3 critique turns, best-of
\newcommand{\dream}{\mbox{\textsc{Dream}}}                 % one dream, up to 3 match turns
\newcommand{\dreamcap}{\mbox{\textsc{Dream-Caption}}}      % Dream + caption of the dream
\newcommand{\shufcap}{\mbox{\textsc{Shuffled-Caption}}}    % Dream + caption of another prompt's dream
\newcommand{\textspec}{\mbox{\textsc{Text-Spec}}}          % spec from the prompt, no image
\newcommand{\bestoffour}{\mbox{\textsc{Best-of-4}}}        % 4 independent no-image builds, pick best
\newcommand{\voxelar}{\mbox{\textsc{Voxel-AR}}}            % 27M trained voxel transformer (recanon20m)
\newcommand{\realset}{\mbox{\textsc{Real}}}                % held-out real houses

\newcommand{\mvkid}{\mbox{MV-KID}}
\newcommand{\geomkid}{\mbox{Geom-KID}}
\newcommand{\blockscore}{\mbox{BlockScore}}
\newcommand{\enterable}{\mbox{enterable}}
\newcommand{\twinrate}{\mbox{twin rate}}

\newcommand{\gptmini}{\mbox{gpt-5-mini}}
\newcommand{\imgmini}{\mbox{gpt-image-1-mini}}


\definecolor{cvprblue}{rgb}{0.21,0.49,0.74}
\usepackage[pagebackref,breaklinks,colorlinks,allcolors=cvprblue]{hyperref}

\def\confName{CVPR}
\def\confYear{2027}
\def\paperID{*****}

%%%%%%%%% TITLE
\title{What Does Visual Imagination Buy a Program-Writing Agent?\\Appearance, Not Structure}

\author{Anonymous CVPR workshop submission}

\begin{document}

\twocolumn[{%
\renewcommand\twocolumn[1][]{#1}%
\maketitle
% Figure 1, full width on page 1.
%TL;DR: one prompt, three arms and the real house show the whole paper: the caption moves the look toward the dream, and no arm builds a house a player walks into as often as real ones.
\begin{center}
  \centering
  \fbox{\parbox[c][4.4cm][c]{0.96\textwidth}{\centering
    \TODO{Figure 1. One row per prompt (two prompts). Columns: (a) the prompt, a houses32-v1 test title;
    (b) the dream drawn by \imgmini{}; (c) the \classicm{} build; (d) the \dream{} build;
    (e) the \dreamcap{} build; (f) the real GrabCraft house with the same title.
    Under (c) to (f) a cutaway render with the player-walk flood fill overlaid, marked \enterable{} yes or no.
    Sources: \texttt{builds/<arm>/NNN/dreams/*.png} and \texttt{iters/} in
    \texttt{outputs/run\_20261007\_092654\_bench\_houses32\_\{dream,classic\}} and
    \texttt{outputs/run\_20261007\_145754\_bench\_houses32\_dream\_caption};
    real houses from the houses32-v1 test split; walk overlay from \texttt{blockgen/eval/traversal.py};
    textured renderer, same camera for every column. Pick prompts present in all three runs.}}}
  \captionof{figure}{\textbf{Imagination buys appearance, not structure.} A language model writes a
  build program for each prompt. With a dream image and a caption of that image (e), the builds
  look closer to real houses in multi-view renders than those of a compute-matched loop without
  an image (c), and the image alone (d) does not move them. In geometry the three arms are not separated, and
  a player enters 50\% to 68\% of the builds, against 93\% of real houses (f).}
  % 0.50 / 0.51 / 0.68 enterable, real 0.95: results.md T63
  \label{fig:teaser}
\end{center}

\vspace{0.5em}
}]

\begin{abstract}
%TL;DR (X) A language model that writes build programs makes Minecraft houses from text for cents, and a natural upgrade is to let it imagine the target as an image first.
%TL;DR (Y) Whether the image helps is hard to measure, because a dream loop also adds turns and text, and its extra text is what the builder reads.
%TL;DR (Z) We separate the image from the turns with a compute-matched loop without an image, and the image from its description with caption controls.
%TL;DR (V) The caption of the dream moves multi-view realism; the image alone does not; no arm moves geometry, and all arms build houses players rarely enter and that repeat each other.
A language model that writes short WorldEdit-style build programs turns a text prompt into a
Minecraft house for about four cents, with no training data. A popular upgrade for such agents
is visual imagination: a text-to-image model draws the target, and the agent builds, renders,
compares and repairs until its build matches the picture. Measuring what the picture adds is
harder than it looks, because an imagination loop also adds repair turns and gives the builder
extra text, and either one alone changes the builds. We isolate the image with a no-image loop
that matches the number of turns and the selection rule, and we isolate the information in the
image with a structured caption of it, a caption of a different prompt's image, and a text
specification written without any image. On 133 house prompts, with \gptmini{} as the builder and
\imgmini{} as the image model, the caption arm nearly halves the multi-view render distance to real
houses against the no-image loop (\mvkid{} 0.217 against 0.403, Holm-corrected $p=0.014$), while
the image without a caption ties it. No arm separates from another on geometry, the
no-image loop builds the most enterable houses (68\% against 50\% and 51\%, real 93\%), and in every
arm a fifth to a third of the builds have a near-copy in the same set, against none of the real houses. Visual imagination
buys a program-writing agent appearance, not structure. \TODO{n=256: replace every number in this
abstract with the n=256 card (lands Nov 10) and add the shuffled-caption, text-spec and best-of-4 results.}
% 0.04 USD/build: card run_20261109_121625 cost table (0.040 to 0.048), T51
% MV-KID 0.217 vs 0.403, p_holm 0.009: T51 addendum 2 (n=133); the 2026-11-09 card gives p_holm 0.0136 with more arms
% enterable 0.68 / 0.50 / 0.51, real 0.95: T63
% twin rate 0.21 to 0.32 vs real 0.008 (ratios 26 to 40): card run_20261109_121625 novelty table, real from T91 (n=256)
\end{abstract}

\section{Introduction}
\label{sec:intro}

%TL;DR (X) Generating buildings in block worlds is a standing 3D generation task, and trained generators need a corpus and a training run for every new condition.
Block worlds such as Minecraft hold millions of player-made buildings, and generating new ones is
a standing task for procedural content and 3D generation~\citep{salge2018gdmc,chen2019voxelcnn,merino2026dreamcubed}.
A building in such a world is a voxel grid in which every cell holds one of hundreds of block
types. Trained generators write this grid cell by cell or denoise it as a
volume~\citep{chen2019voxelcnn,merino2026dreamcubed,pun2025brickgpt}, and each new form of
conditioning, such as a text prompt or a reference image, requires a labelled corpus and a new
training run.

%TL;DR (X) A language model that writes a short build program instead gets text conditioning for free, at a few cents per house.
A language model that writes a build program avoids both. The program is a few dozen commands
in the style of WorldEdit~\citep{worldedit}, such as \texttt{fill}, \texttt{walls} and
\texttt{gable}, and one command places tens to hundreds of blocks. The prompt is the conditioning,
and the model needs no Minecraft training data. With \gptmini{}~\citep{openai2025gpt5} as the
builder, one house on a $32^3$ canvas costs 4.0 to 4.8 US cents in our runs (\cref{sec:method}).
% 0.040 / 0.045 / 0.048 USD per build: card run_20261109_121625 cost table; T51
Several systems already let language models build in Minecraft from text or from a given
reference image~\citep{hu2024t2bm,chen2024apt,fu2026craftutopia,wei2025mineanybuild}.

%TL;DR (X) A popular upgrade for program-writing agents is visual imagination: draw the target image, then build toward it with render-and-compare repair.
A popular upgrade for such agents is visual imagination. A text-to-image model first draws what
the result looks like, and the agent then acts toward that picture. Robot policies
condition on generated subgoal images or videos~\citep{black2024susie,du2023unipi}, web agents
simulate the outcome of an action before taking it~\citep{gu2024webdreamer}, and code-writing
3D agents render their program, compare the render with a target image and edit the
code~\citep{yin2026viga,li2026legoanything}. For a building the appeal is direct: a picture fixes
the massing, the materials and the roof at once, which a short prompt such as
``Small Brick Library (a Minecraft brick house)'' leaves open.

%TL;DR (Y) Measuring what the image adds is hard: an imagination loop also adds repair turns and extra text, and comparing against a one-shot builder credits the image with their effect.
What the image adds is harder to measure than it looks. An imagination loop changes three
things at once. It gives the builder a target picture; it adds repair turns in which the builder
sees its own render; and, in most designs, it adds text about the picture that the builder reads.
A comparison against a builder that writes one program and stops credits the image with the
effect of all three. The image is also an incomplete target for a house: it shows one exterior
view, while what makes a house a house, rooms, doors and stairs a player walks through, sits
behind its walls.

%TL;DR (Z) We build a recipe harness in which every arm differs from its control in one ingredient, with images shared across arms through a content-addressed cache.
We therefore run imagination as a controlled study. A recipe harness defines each arm as one
YAML file over the same builder, build language, renderer and budget, and a content-addressed
image cache gives every arm that draws a dream exactly the same images (\cref{sec:method}).
\classicm{} runs the same number of repair turns and the same best-of selection as the
imagination loop, but critiques its own renders without a target image. \dream{} adds the target
image. \dreamcap{} adds a structured caption of that image, written by a captioner that never sees
the prompt. Three controls separate the caption's content from its length and from the image
itself: \shufcap{} shows the caption of another prompt's dream, \textspec{} shows a specification
written from the prompt with no image, and \bestoffour{} spends the same compute on independent
samples without any repair.

%TL;DR (1a) Finding on appearance: the caption nearly halves the multi-view render distance to real houses, and the image without the caption ties the no-image loop.
On 133 house prompts, the caption carries the visual gain. \dreamcap{} reaches a multi-view
render distance to real houses (\mvkid{}) of 0.217, against 0.403 for \classicm{} and 0.342 for
\dream{}, and beats both in Holm-corrected paired tests ($p=0.014$). \dream{} and \classicm{} share
a significance group. The image alone therefore buys nothing a matched loop without an image
does not, and the gain appears once the image is turned into words the builder reads.
\TODO{n=256 numbers; \shufcap{} and \textspec{} decide whether the gain is the caption's content or its length.}
% MV-KID 0.217 / 0.342 / 0.403, p_holm 0.009 both pairs, dream~classic same group: T51 addendum 2; card run_20261109_121625 (groups d / e / ef)

%TL;DR (1b) Finding on structure: no arm moves geometry, the no-image loop builds the most enterable houses, and all arms repeat themselves far more than real houses do.
Structure does not follow appearance. On the voxel-geometry distance (\geomkid{}) no agentic arm
separates from another (\dream{} 17.9, \classicm{} 18.4, \dreamcap{} 15.7), and a 27M-parameter
voxel transformer trained on houses (\voxelar{}) reaches 9.06. A simulated player enters 68\% of
\classicm{} builds, 50\% of \dream{} builds and 51\% of \dreamcap{} builds, against 93\% of real
houses. Within each arm, 21\% to 32\% of builds have a near-copy in the same set, against none
of the real houses, so all three arms fail the near-twin gate of the evaluation protocol we
adopt~\citep{anonymous2027blindness}. None of them copies a training house.
\TODO{n=256: replace \geomkid{}, \enterable{} and \twinrate{} with the n=256 card; check whether the
distinct-prompt card lowers the twin rate.}
% geom_kid 17.882 / 18.351 / 15.712 (groups ef / ef / de), recanon20m_s0 9.063: card run_20261109_121625; 27M params: T118
% enterable 0.68 / 0.50 / 0.51, real 0.93: card run_20261109_121625 (T63 gives 0.95 at n=128)
% twin rate 0.211 / 0.301 / 0.316: card run_20261109_121625; real 0.000 on this card (0.008 at n=256, T91); voxel_dup_rate 0 for all agentic arms: same card

%TL;DR (C) Contributions: a one-ingredient-at-a-time study of imagination, the finding that captions carry the visual gain, and the finding that neither image nor caption reaches structure.
Our contributions are:
\begin{itemize}[leftmargin=*,itemsep=1pt,topsep=2pt]
  \item A controlled study of visual imagination for a program-writing building agent, in which
    each arm differs from its control in one ingredient: turns, image, caption, caption
    content, or sampling (\cref{sec:method}).
  \item Evidence that the visual gain comes from a caption of the imagined image and not from the
    image itself: the caption nearly halves the multi-view render distance to real houses, and the image
    alone ties a compute-matched loop without an image (\cref{sec:results}).
  \item Evidence that neither the image nor the caption reaches structure: geometry is not
    separated, the loop without an image builds the most enterable houses, and every arm repeats
    near-identical builds far more often than real houses do (\cref{sec:results}).
\end{itemize}

\section{Related Work}
\label{sec:related}

%TL;DR Trained block-world generators learn structure from a corpus; we test a training-free agent against them and find the trained model still ahead on geometry.
\paragraph{Trained generators for block worlds.}
Learned generators for Minecraft write a building block by block in the order a player placed
it~\citep{chen2019voxelcnn}, denoise whole chunks after training on billions of
cubes~\citep{merino2026dreamcubed}, or optimize a block field against a text-image score for each
prompt~\citep{earle2024dreamcraft}. For bricks, \citet{pun2025brickgpt} train an autoregressive
model on a text-annotated brick corpus and reject unstable placements during decoding. All of these
learn structure from data or from per-prompt optimization, and each new kind of condition needs a
new corpus or objective. Our agent needs neither, which makes it cheap to condition on an imagined
image; we keep a trained voxel transformer in the comparison to show what the agent gives up in
geometry.

%TL;DR LLM building agents in Minecraft take text, a plan or a given reference image; none measures whether a self-drawn image helps against a loop matched in turns.
\paragraph{Language-model building agents in Minecraft.}
\citet{hu2024t2bm} prompt language models to write Minecraft buildings from text, and
\citet{chen2024apt} decompose the prompt into an architectural plan and blueprint with memory and
reflection modules. CraftUtopia~\citep{fu2026craftutopia} turns a given 2D reference image into a
blueprint and splits construction across parallel workers. Word2Minecraft~\citep{huang2025word2minecraft}
generates game levels from text, Voyager~\citep{wang2023voyager} writes skill code to act in the game, and
MineAnyBuild~\citep{wei2025mineanybuild} and MC-Bench~\citep{mcbench2025} evaluate agents and
models that build from instructions. In these systems the image, when present, is an input the user
provides. We ask a different question: whether an image the agent draws for itself improves the
build, measured against a loop that spends the same turns without one.

%TL;DR Code-writing 3D agents with render feedback chase a given target image or a validator; our target is self-made, so the open question is whether it beats having no target.
\paragraph{Code-writing 3D agents with visual feedback.}
Agents that write graphics programs and repair them from renders are common outside block worlds.
SceneCraft~\citep{hu2024scenecraft} and LL3M~\citep{lu2025ll3m} write Blender code for scenes and
assets, CADCodeVerify~\citep{alrashedy2025cadcodeverify} checks CAD code against renders of its
output, and VIGA~\citep{yin2026viga} and LEGO-Anything~\citep{li2026legoanything} run a
write-render-compare loop until the program reproduces a target image. For bricks,
BrickWright~\citep{pun2026brickwright} and BrickBench~\citep{kulits2026brickbench} give coding
agents a brick assembly API with a physical validator. In inverse graphics the target is given and
the render-to-target distance is the task; in our setting the target is itself generated, so a
build that matches it perfectly is only as good as the picture. Benchmarks of code-written 3D
report that executable code is often not spatially correct~\citep{zheng2026voxelcodebench} and that
successful renders often hide floating or disconnected geometry~\citep{gao2026threedcodebench}. We
find the analogous gap one level up: a build that looks closer to real houses in renders is not
closer in geometry or in walkability.

%TL;DR Imagined goals help robot and web agents because there the image fixes the goal state; one exterior view underspecifies a building, which predicts our result.
\paragraph{Imagined goals for acting.}
Generating the goal before acting helps embodied agents: SuSIE~\citep{black2024susie} edits the
current camera image into a subgoal for a low-level policy, UniPi~\citep{du2023unipi} plans as a
generated video and reads actions from it, and WebDreamer~\citep{gu2024webdreamer} uses a language
model to simulate the outcome of each web action. CaP-X~\citep{fu2026capx} studies code-writing robot
agents and lists visual differencing among the test-time techniques that narrow the gap to human-designed
abstractions. In manipulation a subgoal image specifies most of the state that matters. A single
exterior view of a house does not: it fixes the facade and leaves the interior, the back and the
circulation open. Our results match this difference: the imagined image moves appearance and leaves
structure where it was.

%TL;DR We adopt an existing evaluation protocol unchanged and make no claim about evaluation.
\paragraph{Evaluation.}
We score every arm with the protocol of a concurrent submission~\citep{anonymous2027blindness}
without modification, including its distribution metrics, its walkability probe and its novelty
gates, and we make no claim about the protocol itself.

\section{Arms and Controls}
\label{sec:method}

%TL;DR Every arm shares one builder, one build language, one renderer and one canvas, and differs from its control in exactly one ingredient.
All arms share one builder model (\gptmini{}), one build language, one textured renderer and one
$32^3$ canvas, and each arm is a recipe file that differs from its control in one ingredient.
\Cref{tab:arms} lists the ingredients. This section describes the build language
(\cref{sec:language}), the loop without an image (\cref{sec:classic}), the imagination loop
(\cref{sec:dream}), the caption (\cref{sec:caption}), the controls (\cref{sec:controls}) and the
measures that keep the comparison fair (\cref{sec:fair}).

\begin{table}[t]
  \centering
  \small
  \setlength{\tabcolsep}{3pt}
  \begin{tabular}{@{}lcccr@{}}
    \toprule
    Arm & Image & Target text & Turns & \$/build \\
    \midrule
    \classicm{}  & no  & none                        & $1{+}3$        & 0.045 \\
    \dream{}     & yes & none                        & $1{+}{\le}3$   & 0.040 \\
    \dreamcap{}  & yes & own caption & $1{+}{\le}3$   & 0.048 \\
    \midrule
    \shufcap{}   & yes & other caption & $1{+}{\le}3$   & \tbd{} \\
    \textspec{}  & no  & spec & $1{+}3$        & \tbd{} \\
    \bestoffour{}& no  & none                        & $4{\times}1$   & \tbd{} \\
    \bottomrule
  \end{tabular}
  % $/build: card run_20261109_121625 cost table (cost_usd, n=133); T51 reports the same 0.045 / 0.040 / 0.048
  \caption{\textbf{Arms.} ``Turns'' counts one repair turn for execution errors plus the critique
  or match turns; \bestoffour{} draws four independent builds with one repair turn each. All arms
  select the best attempt they produce. The lower three rows are controls; their cost is
  \TODO{n=256 card}.}
  \label{tab:arms}
\end{table}

\subsection{Build language}
\label{sec:language}

%TL;DR The builder writes a short program in a WorldEdit-style language whose commands place many blocks each, over a fixed palette of about seventy block names.
The builder writes a program of one command per line in a language modelled on the in-game
editor WorldEdit~\citep{worldedit}. Commands take inclusive integer coordinates with $y$ up and
the origin at a canvas corner: \texttt{set} places one block, \texttt{fill}, \texttt{clear},
\texttt{box}, \texttt{walls} and \texttt{replace} act on cuboids, \texttt{line},
\texttt{sphere}, \texttt{cylinder}, \texttt{pyramid} and \texttt{gable} draw primitives, and
\texttt{copy}, \texttt{paste} and \texttt{stack} repeat parts. Blocks are named from a curated
palette of about seventy modern names, with block states such as stair orientation where the game
stores them. The system prompt lists the commands and the palette, gives one example program, and
states player scale: a player is two blocks tall and one block wide, so doors are two blocks tall.

%TL;DR Parsing and execution report errors per line instead of failing, which gives every repair turn a precise list of what to fix.
The executor parses the whole program before running it and reports syntax errors with line
numbers. A command that fails is recorded and skipped while the rest executes, writes outside the
canvas are clipped and counted, an unknown block name is an error on its line rather than a
silent substitution, and a command that changes no block raises a warning, which almost always
marks a coordinate mistake. Every arm spends one repair turn on these reports after its first
program.

\subsection{The loop without an image: \classicm{}}
\label{sec:classic}

%TL;DR The compute-matched control writes a program, repairs execution errors, then critiques and rewrites from two renders of its own build three times, and keeps the best attempt.
\classicm{} is the control for every arm that draws an image. After the first program and its
repair turn, it runs three critique turns. In each, the builder sees renders of its current build
from two fixed cameras and rewrites the whole program. A new attempt is accepted when its score
improves and the share of blocks in its largest connected component drops by at most 0.1, which
blocks the cheap trick of adding floating decoration. The arm returns its best attempt under a
score that does not depend on any image target: the similarity between the prompt and the render
under CLIP ViT-B/16~\citep{radford2021clip}, plus half the largest-component share.

\subsection{The imagination loop: \dream{}}
\label{sec:dream}

%TL;DR The dream loop draws one isometric target image, builds toward it, and then compares render and target in up to three match turns.
\dream{} adds a target image, the dream, and nothing else (\cref{alg:dream}). The image model
(\imgmini{}~\citep{openai2025gptimage}, low quality, $1024{\times}1024$) draws the prompt from a
fixed isometric camera in a pinned style: a single build of full cube blocks in flat daylight on a
white background, with no terrain, sky, people or text. The template also states the canvas size
and player scale. The builder sees the dream labelled as the target in its first turn. After the
repair turn, up to three match turns follow. In each, the builder sees the target next to a render
of its build from the dream's camera, labelled as its own, and rewrites the program.

\begin{algorithm}[t]
  \caption{\dream{}, \dreamcap{} and \shufcap{}. Lines marked $\dagger$ run only in the caption arms.}
  \label{alg:dream}
  \small
  \begin{algorithmic}[1]
    \Require prompt $p$, canvas $32^3$, at most $K{=}3$ match turns
    \State $I \gets \textsc{ImageModel}(\text{style}, p)$ \Comment{cached by request}
    \State $c \gets \textsc{Captioner}(I)$, or another prompt's caption \Comment{$\dagger$}
    \State $P_0 \gets \textsc{Builder}(p, I, [c])$; repair once from executor errors
    \For{$k = 1, \dots, K$}
      \State \textbf{if} $s(P_{k-1}) \ge 0.85$ or two turns without gain \textbf{then break}
      \State $R \gets$ render of $P_{k-1}$ at the camera that best matches $I$
      \State $P_k \gets \textsc{Builder}(\text{target } I, \text{yours } R, [c])$
      \State accept $P_k$ if $s$ improves and connectivity holds
    \EndFor
    \State \Return $\arg\max_k s(P_k)$
  \end{algorithmic}
\end{algorithm}

%TL;DR The in-loop score compares render and dream with backbones disjoint from the evaluation's, gated on connectivity, and the loop stops early on success or stall.
The loop scores each attempt against the dream with the cosine similarity of SigLIP-B
embeddings~\citep{zhai2023siglip} (weight 1.0), a silhouette overlap (weight 0.5) and the
largest-component share (weight 0.3); an attempt whose largest component holds less than half its
blocks scores zero. The render is taken from whichever of the four isometric cameras matches the
dream best, because the image model's camera is approximate. The loop stops when the score reaches
0.85 or after two turns without improvement, applies the same acceptance rule as \classicm{}, and
returns the best attempt. Neither in-loop backbone (SigLIP-B here, CLIP ViT-B/16 in \classicm{})
appears among the evaluation's backbones~\citep{anonymous2027blindness}, so no arm optimizes the
quantity it is ranked by.

\subsection{Describing the dream: \dreamcap{}}
\label{sec:caption}

%TL;DR The caption arm adds one vision turn in which a captioner, who never sees the prompt, describes the dream in a fixed schema that the builder reads in every turn.
\dreamcap{} adds one vision turn to \dream{} and changes nothing else. A captioner (also
\gptmini{}) receives the dream and the canvas size, but not the prompt, and describes what the
image shows in a fixed JSON schema: a one-sentence summary, the size and footprint in blocks, the
storey count, the roof shape and material, the palette blocks of the walls, trim, roof, windows,
door and foundation, the openings, distinctive features, notes on the front and side faces,
anything around the building to leave out, and a list of what is uncertain. Fields the image does
not show, such as the back or the interior, are marked unknown. The builder reads the caption as
a description of the target in its first turn and in every match turn. Because the image cache
returns the same dream for the same request, \dreamcap{} and \dream{} build toward identical
images, and any difference between them is the caption.

\subsection{Controls}
\label{sec:controls}

%TL;DR Three controls ask whether the caption helps through its content, whether the image is needed for that content, and whether more samples alone do as well.
Each control answers one question about the caption result.
\begin{itemize}[leftmargin=*,itemsep=1pt,topsep=2pt]
  \item \shufcap{} makes the same calls as \dreamcap{} but shows the builder the caption of the
    dream of the previous prompt in the same arm; the first prompt receives the one of five fixed
    library captions that shares the fewest words with it. Since the captioner never sees the
    prompt, a shown caption carries no trace of where it came from. If \shufcap{} matches
    \dreamcap{}, the gain is extra text, not description.
  \item \textspec{} writes a specification of the house from the prompt alone, in the fields
    of the caption without the size and per-face notes, plus the floor material and storey heights, and runs the \classicm{}
    loop with the specification in place of an image. If \textspec{} matches \dreamcap{}, the
    image is not needed: a structured description is. \TODO{\textspec{} has no recipe yet in
    \texttt{configs/agentic/recipes}; the closest is \texttt{dream-describe-before}, which draws
    an image from the spec. Add a no-image recipe and confirm this description matches it.}
  \item \bestoffour{} draws four independent \classicm{} builds with no critique turns and keeps
    the best under the \classicm{} score. If it matches \dream{} or \dreamcap{}, the gain is
    sampling, not imagination.
\end{itemize}

\subsection{Fair comparison}
\label{sec:fair}

%TL;DR Shared caches, recipe hashes and a cost ledger make arms paired on images, exactly reproducible, and comparable in dollars.
Three mechanisms keep the arms comparable. First, dreams and model replies sit in
content-addressed caches, so arms that issue the same request receive the same response, and a
rerun in read-only mode reproduces a whole arm at zero cost. Second, each recipe resolves to a
hash over all of its fields, and two runs count as the same arm only when their hashes match.
Third, a ledger prices every call before it is made, enforces a budget per build and per run, and
records both the price paid and the full price of calls served from the cache.
All arms run on the same prompts, the titles of held-out houses written as
``\textit{title} (a Minecraft \textit{category})'', so each build has a real house with the same
name in the evaluation set (\cref{sec:setup}).
\TODO{n=256: confirm the 123 added prompts are validation-split titles in the same format, and that the card reports full-price cost.}

\section{Experimental setup}
\label{sec:setup}

%TL;DR Prompts are the GrabCraft titles of the held-out test houses, so every prompt has a real house it names and the generated set has a real distribution to match.
\paragraph{Prompts.}
Each prompt is the GrabCraft title of a held-out house, phrased as ``\texttt{<title> (a Minecraft <category>)}'', for example ``Viking Inn (a Minecraft medieval house)''. % T51; configs/agentic/prompts/houses32_test_grabcraft.txt
The main comparison uses the 133 test houses of the \texttt{houses32-v1} split that carry a GrabCraft title. % T51: 133 test titles
Their categories are medieval (84), modern (19), wooden (10), survival (6), stone (4), tree (4), starter (3), brick (2) and quartz (1). % counted from houses32_test_grabcraft.txt
To reach the protocol's headline size of $n{=}256$, we add 123 titles of validation houses, disjoint from the test titles. % configs/agentic/prompts/houses32_val_grabcraft.txt; val run config: 123 prompts; file holds 124 distinct titles, 0 shared with test
\TODO{state why the val run uses 123 of the 124 distinct titles in houses32\_val\_grabcraft.txt (the dropped title is ``Medieval City Hall With a Bell Tower (a Minecraft wooden house)'').}
Every arm receives the same prompts in the same order, and the builder never sees a real house.

%TL;DR All arms share one builder (gpt-5-mini at medium reasoning effort), one action space (WorldEdit-style commands; 32^3 canvas for the loop arms, 48^3 for One-Shot) and one textured renderer; the dream arms add gpt-image-1-mini at low quality.
\paragraph{Models.}
Every text and vision role (builder, critic, judge, captioner) is \texttt{gpt-5-mini}~\citep{openai2025gpt5} with reasoning effort ``medium'', default temperature and a 16{,}000-token completion cap. % configs/agentic/recipes/_base.yaml
The builder writes a WorldEdit-style program that an executor runs on a $32^3$ canvas, and a textured renderer draws the result for the critic. % _base.yaml: actions dsl, canvas [32,32,32], renderer textured
The dream arms generate one $1024{\times}1024$ isometric target per build with \texttt{gpt-image-1-mini}~\citep{openai2025gptimage} at low quality. % dream.yaml: image openai:gpt-image-1-mini?quality=low&size=1024x1024; n_images 1.000 per build
\textsc{One-Shot} runs on a $48^3$ canvas and its builds are reduced to the protocol's 32-voxel maximum by the benchmark's structure cache. % oneshot config.json: canvas [48,48,48]; output legacy-oneshot_32.npz
\TODO{confirm how the $48^3$ One-Shot builds reach max\_dim 32 (crop or rescale) and whether to rerun One-Shot at $32^3$ for parity.}

%TL;DR Four arms are measured: One-Shot, Classic-Matched, Dream and Dream-Caption; each differs from its neighbour by one component.
\paragraph{Arms.}
\textsc{One-Shot} writes one program from the prompt and one in-context example, with no repair and no critique. % oneshot config: plan False, n_examples 1, repair_rounds 0, critique_rounds 0
\textsc{Classic-Matched} adds one executor-error repair turn and three render-critique turns, each of which rewrites the program after seeing a render of the current build. % classic-matched.yaml: repair_rounds 1, critique_rounds 3, critique_mode rewrite
\textsc{Dream} first draws a target image, builds from it, renders the build from the target's camera, and runs up to three compare-and-rewrite iterations (patience 2, early stop at an in-loop score of 0.85). % dream.yaml: max_iters 3, patience 2, stop_at 0.85
\textsc{Dream-Caption} adds one vision turn in which a captioner that never sees the prompt describes the image (materials as palette blocks, storeys, roof, dimensions, per-face notes, an uncertainty list); the builder and the critic receive this caption beside the image. % dream-caption.yaml: describe after; docs/agentic-suite.md 5.6
\textsc{Dream} and \textsc{Dream-Caption} share an image cache, so both arms build from the same target image for every prompt and their difference is the caption alone. % dream-caption.yaml comment
\textsc{Classic-Matched} and both dream arms carry the same player-scale rules in their context. % classic-matched.yaml / dream.yaml: contexts include scale

%TL;DR Classic-Matched is the compute-matched control: it receives at least as many builder turns as Dream, and its in-loop scorer shares no backbone with the benchmark.
\paragraph{Compute matching.}
\textsc{Classic-Matched} exists to isolate the image: any gain of \textsc{Dream} over it is attributable to the target image and not to more turns. % classic-matched.yaml comment
Measured over the 133 test prompts, \textsc{Classic-Matched} makes 4.03 model calls per build against 3.49 for \textsc{Dream} and 4.56 for \textsc{Dream-Caption}, and costs \$0.045 per build against \$0.040 and \$0.048. % run_20261007_092654_bench_houses32_{classic,dream}/summary.md, run_20261008_004051_bench_houses32_dream_caption133/summary.md: n_llm_calls, would_cost/133; scorecard cost table
\textsc{Dream-Caption} therefore spends 7\% more than the control, one extra call per build. % 6.385/5.968 = 1.07, would_cost totals
Each arm selects its returned build by an in-loop score: \textsc{Classic-Matched} by CLIP text similarity~\citep{radford2021clip} plus a validity term, the dream arms by SigLIP~\citep{zhai2023siglip} image similarity to the target plus silhouette and validity terms. % classic-matched.yaml score {clip_text 1.0, validity 0.5}; dream.yaml score {embed 1.0, silhouette 0.5, validity 0.3}; docs/agentic-suite.md 9: in-loop siglipB / clipB16
Both in-loop backbones are disjoint from the DINOv2 and CLIP-L backbones of the benchmark, so no arm optimizes the metric it is scored on. % docs/agentic-suite.md section 9 (Goodhart)

%TL;DR Planned controls fill the ablation table at n=256: a shuffled caption, a text-only spec, best-of-4 sampling, symbolic inspection feedback and a local model.
\paragraph{Planned controls.}
\textsc{Dream-Caption-Shuffled} makes the same calls as \textsc{Dream-Caption} but shows the builder the caption of a different prompt's image, which separates a correct description from extra text. % dream-caption-shuffled.yaml
\textsc{Text-Spec} gives the builder a written design spec and no image, which separates the description from the picture. \TODO{no recipe named text-spec exists in configs/agentic/recipes; define it (closest existing: dream-describe-before without the image) before the run.}
\textsc{Best-of-4} keeps the best of four independent builds under \textsc{Classic-Matched}'s selector, which separates the dream from sampling more. % best-of-4.yaml
Two further arms are built but unmeasured: \texttt{inspect-lessons}, whose critic receives text geometry checks (grounded, connected, enterable, headroom) and retrieved building lessons in place of a render, and a local Gemma~4 12B builder with SDXL-Turbo dreams. % inspect-lessons.yaml, docs/agentic-suite.md 5.7; dream-caption-gemma-local.yaml, T48

%TL;DR We score every arm with the protocol of a concurrent submission: Geom-KID and MV-KID as co-primary distances, plus a traversal probe, all against 399 held-out real houses.
\paragraph{Evaluation.}
We use the evaluation protocol of a concurrent submission~\citep{anonymous2027blindness}, unchanged. % scorecard: reference 399, seed 0
MV-KID is a kernel inception distance~\citep{binkowski2018kid} on DINOv2~\citep{oquab2024dinov2} features of multi-view textured renders and reads appearance. % scorecard realism table: mv_dino_kid
Geom-KID is the same distance on an 80-dimensional voxel descriptor of occupancy patterns, thickness, height and the player's room scalars, and reads geometry. % T64: 80-dim descriptor with room scalars
The protocol reports the two as co-primary, because render metrics miss interiors and structure. % anonymous2027blindness; memory eval-suite-bench
A traversal probe sets a two-block player down outside each build and flood-fills every cell it walks to; a build is \emph{enterable} when the player reaches an interior standing cell. % T63
The reference is 399 real test houses, and a held-out real sample scored as an arm gives the noise floor. % scorecard: reference 399, real_test arm

%TL;DR Gates disqualify an arm for training-set memorization, for mode collapse, and for near-twins inside its own set, before any score is read.
\paragraph{Gates.}
Three gates act before ranking. The memorization gate disqualifies an arm whose builds sit closer to the training set in DINOv2 space than held-out real builds do (DINO-NN percentile). % scorecard leaderboard: dino_nn_percentile memorization gate
The diversity gate disqualifies mode collapse. % scorecard: voxel_diversity gate (real@single_mode)
The near-twin gate counts the share of builds that have a near-copy inside the same set (occupancy IoU at least 0.8 under the eight footprint symmetries and translation) and disqualifies an arm more than three real spreads above the real rate, with the spread floored at 0.02. % T91
An arm that passes all three receives a BlockScore, the worst of its realism pillars in units of the real-versus-real spread. % scorecard footnote; docs/benchmark.md

%TL;DR We test every pair of arms with paired permutation tests under Holm correction, and we compare BlockScore only between arms scored at the same n.
\paragraph{Statistics.}
We compare arms with permutation tests (10{,}000 permutations), scoring both arms of a pair against the same reference, and correct the whole family of pairs on the card with Holm's method~\citep{holm1979simple}. % scorecard: "paired, Holm-corrected ... permutation, n_perm=10000"; docs/benchmark.md head-to-head
We report the resulting compact letter display: arms that share a letter are not separated at $\alpha{=}0.05$. % scorecard
The Holm family contains every pair among 17 arms of the card, so it is conservative for the four agentic arms. % scorecard arms: 17 in mv_dino_kid ranking
Intervals are 95\% bootstrap intervals over builds. % scorecard: bootstrap unit structure
BlockScore is calibrated against a real sample of the arm's own $n$, so we compare it only between arms at equal $n$. % docs/benchmark.md Protocols
The calibration matters in practice: the same \texttt{canon\_oriented} samples score BlockScore 150.5 on the $n{=}133$ card and 297 on an $n{=}256$ card. % scorecard (150.53, real_test n=133); T118 (297)
Every arm runs once, with one build per prompt, and the benchmark uses seed 0. % scorecard: seed 0; summaries: n_builds 133
All agentic numbers in \cref{sec:results} come from one card scored at $n{=}133$. % outputs/run_20261109_121625_bench_agentic_arms_beside_the_20M_legacy/scorecard.md
\TODO{n=256: replace with the n=256 card once the val runs of Dream-Caption, Dream and Classic-Matched finish and One-Shot is benched.}

\section{Results and discussion}
\label{sec:results}

% Source of every agentic and trained-model number in this section unless noted:
% outputs/run_20261109_121625_bench_agentic_arms_beside_the_20M_legacy/scorecard.md
% (houses_32.s0.70-15-15.v1, seed 0, reference 399, 17 arms, permutation n_perm=10000, Holm).
% p_Holm = 0.0136 is the smallest Holm-adjusted p the card reports; we print it as 0.014.

\begin{table*}[t]
\centering
\small
\setlength{\tabcolsep}{4pt}
\caption{\textbf{Agentic arms on the house benchmark, $n{=}133$ test prompts.} MV-KID reads appearance and Geom-KID geometry (lower is better; letters: arms sharing a letter are not separated after Holm correction). Enterable is the share of builds a player walks into (real reference 0.945). Twin is the share of builds with a near-copy in the same set. Trained models are scored on the same card at $n{=}256$, so their BlockScore is not comparable with an $n{=}133$ arm (\cref{sec:setup}). Brackets are 95\% intervals.}
\label{tab:main}
\begin{tabular}{@{}llcccccl r@{}}
\toprule
Arm & $n$ & MV-KID $\downarrow$ & Geom-KID $\downarrow$ & Enterable & Twin & Diversity & Gate & \$/build \\
\midrule
Real (held out) & 133 & $-0.001$\,\textsuperscript{a} & $-0.08$\,\textsuperscript{a} & 0.932 & 0.000 & 0.902 & pass & -- \\ % scorecard real_test rows
\midrule
\textsc{One-Shot} & 133 & \multicolumn{6}{c}{\TODO{bench One-Shot}} & 0.010 \\ % run_20261109_130512_bench_houses32_oneshot_test: paid $1.3892 / 133 builds; not yet benched
\textsc{Classic-Matched} & 133 & 0.403\,\textsuperscript{ef} & 18.35\,\textsuperscript{ef} & 0.677 & 0.316 & 0.833 & DQ twin & 0.045 \\ % scorecard realism/traversal/novelty/cost/leaderboard
\textsc{Dream} & 133 & 0.342\,\textsuperscript{e} & 17.88\,\textsuperscript{ef} & 0.496 & 0.301 & 0.786 & DQ twin & 0.040 \\ % scorecard
\textsc{Dream-Caption} & 133 & \textbf{0.217}\,\textsuperscript{d} & \textbf{15.71}\,\textsuperscript{de} & 0.511 & \textbf{0.211} & 0.790 & DQ twin & 0.048 \\ % scorecard
\midrule
\multicolumn{9}{@{}l}{\emph{Trained voxel models (reference rows)}} \\
\texttt{canon\_oriented} 5M & 256 & 0.254\,\textsuperscript{d} & 14.24\,\textsuperscript{d} & 0.605 & 0.000 & 0.885 & pass (150.5) & -- \\ % scorecard canon_oriented rows; BlockScore 150.53
20M legacy & 256 & 0.159\,\textsuperscript{c} & 9.06\,\textsuperscript{c} & 0.629 & 0.008 & 0.878 & DQ mem. & -- \\ % scorecard recanon20m_s0 rows; T118 model
\bottomrule
\end{tabular}
\end{table*}
% Intervals for the agentic rows (scorecard realism table):
% MV-KID: classic 0.403 [0.361, 0.450], dream 0.342 [0.311, 0.375], caption 0.217 [0.200, 0.237]
% Geom-KID: classic 18.351 [17.056, 19.534], dream 17.882 [17.028, 18.980], caption 15.712 [14.888, 16.681]
% Twin: classic 0.316 [0.241, 0.391], dream 0.301 [0.226, 0.383], caption 0.211 [0.143, 0.286]

%TL;DR We organize the results as six findings, each tied to one row group of the main table; the dream buys appearance through its caption and buys nothing in geometry or structure.
\Cref{tab:main} holds the main comparison. Each subsection below states one finding and the test behind it.

\subsection{The caption, not the image, moves appearance}

%TL;DR Dream-Caption is the only agentic arm separated from the compute-matched control on MV-KID; Dream alone is not.
\textsc{Dream-Caption} reaches MV-KID 0.217, against 0.342 for \textsc{Dream} and 0.403 for \textsc{Classic-Matched}. % scorecard realism
It beats \textsc{Dream} by 0.125 (95\% interval $[0.089, 0.164]$, $p_{\text{Holm}}{=}0.014$) and \textsc{Classic-Matched} by 0.186 ($[0.135, 0.229]$, $p_{\text{Holm}}{=}0.014$). % scorecard mv_dino_kid h2h: dream-caption vs dream -0.125, vs classic-matched -0.186
\textsc{Dream} and \textsc{Classic-Matched} share a letter, so the image alone does not separate from the compute-matched control. % scorecard groups e / ef
Because the two dream arms build from identical cached images, the 0.125 gap belongs to the caption. % dream-caption.yaml shared cache
\textsc{Dream-Caption} also shares its letter with the trained \texttt{canon\_oriented} model (0.254), so on appearance a prompted agent with a captioned image is not separated from a 5M model trained on the corpus. % scorecard groups d / d

%TL;DR Small pilots show what the caption does: it names parts the builder misreads in the image and lists scenery to leave out.
Pilots show the mechanism, at sizes too small to test. % T46 n=4; T48 n=6
In a four-prompt pilot, the caption arm was the only arm whose windmill ended with upright sails, after the caption named ``four large diagonal sail arms''. % T46
In a six-prompt local pilot with SDXL-Turbo dreams that place every house on terrain, the caption listed the terrain as surroundings to omit and the builds dropped it (coherence rate 1.00 against 0.83 without the caption). % T48: coherence 1.00 vs 0.83
In the same four-prompt pilot, a caption from a different prompt turned the lighthouse into a red box house. % T46 dream-caption-shuffled
\textsc{Dream-Caption-Shuffled} tests this at full $n$ (\cref{tab:controls}).

%TL;DR The caption arm builds smaller houses than real ones, which is part of how its renders differ.
The caption also shrinks the builds: \textsc{Dream-Caption} places 704 blocks with a largest extent of 13.4, against 953 and 15.0 for \textsc{Dream}, 1{,}202 and 19.1 for \textsc{Classic-Matched}, and 904 and 19.2 for real houses. % scorecard coherence table: n_blocks, max_dim
\TODO{check whether MV-KID renders normalize build size; if not, measure how much of the caption's MV-KID gain survives a size-matched comparison.}

\subsection{Geometry: the three agentic arms tie}

%TL;DR No pair of agentic arms separates on Geom-KID, so neither the image nor the caption moves the distribution of shapes.
On Geom-KID the three arms score 15.71, 17.88 and 18.35, and every pair shares a letter. % scorecard geom_kid groups de / ef / ef
The image and the caption leave the shape distribution where the compute-matched control puts it.
\textsc{Dream} and \textsc{Classic-Matched} also share a letter with the protocol's Patchwork baseline (18.79), and only \textsc{Dream-Caption} separates from it ($-3.08$, $[-4.52, -1.55]$, $p_{\text{Holm}}{=}0.014$). % scorecard: patchwork 18.791 group f; patchwork vs dream-caption +3.080
All three beat the solidified-real control (21.56), whose interiors are filled. % scorecard: real@solidify 21.562 group g; vs classic-matched p_holm 0.0425

\subsection{Structure: the dream arms are harder to enter}

%TL;DR A player walks into 68% of Classic-Matched houses but only half of the dream houses, against 93% of real ones.
The traversal probe orders the arms the other way from MV-KID. % scorecard traversal table
A player enters 0.677 of \textsc{Classic-Matched} builds, 0.511 of \textsc{Dream-Caption} builds and 0.496 of \textsc{Dream} builds, against 0.932 for held-out real houses. % scorecard enterable
In real-sample spreads the distance to real is 1.17 for the control and 1.90 and 1.96 for the dream arms. % scorecard traversal: enterable spreads
Players reach 0.77 storeys per \textsc{Classic-Matched} build and 0.51 to 0.53 per dream build, against 1.57 for real houses. % scorecard floors_reached
The protocol reports traversal as a distance to real without a pairwise test. \TODO{run a paired permutation test on per-build enterability between Classic-Matched and each dream arm.}

\begin{table}[t]
\centering
\footnotesize
\setlength{\tabcolsep}{3pt}
\caption{\textbf{Inspection failures.} Share of builds failing each text geometry check at warning severity or worse; checks with no failures in any row are omitted. Real houses fail mostly through crop artefacts (missing floors, detached yard pieces).}
\label{tab:inspect}
\begin{tabular}{@{}lrrrrr@{}}
\toprule
Arm & $n$ & Clean & Roofed & Enterable & Headroom \\
\midrule
\textsc{Classic-Matched} & 133 & 68\% & 6\% & 33\% & 1\% \\
\textsc{Dream} & 133 & 53\% & 2\% & 54\% & 7\% \\
\textsc{Dream-Caption} & 133 & 48\% & 0\% & 57\% & 13\% \\
Real houses & 400 & 55\% & 10\% & 17\% & 9\% \\
\bottomrule
\end{tabular}
\end{table}
% Source: T95 (blockgen/agentic/inspection.py on the T51 builds and 400 real houses_32 builds, seed 1).
% Real also fails connected 6%, supported 14%; agentic 0% on both; dream-caption interior 2%.

%TL;DR An independent rule-based inspector agrees: the enterable check fails on more than half of the dream builds and a third of the control builds.
A rule-based inspector, run on the same builds, agrees (\cref{tab:inspect}). % T95
The enterable check fails on 54\% of \textsc{Dream} builds and 57\% of \textsc{Dream-Caption} builds, against 33\% for \textsc{Classic-Matched} and 17\% for real houses. % T95
Headroom failures rise from 1\% for the control to 7\% and 13\% for the dream arms. % T95
The agentic builds never fail the connected or supported checks, which real crops fail at 6\% and 14\%. % T95
The dream arms build connected, supported shells that a player does not get into.

%TL;DR All agentic arms leave interiors too small and too symmetric, and the image moves neither.
Interiors are the shared gap. % scorecard coherence / geometry scalars
The enclosed-air ratio is 0.039 to 0.043 for the three arms against 0.114 for real houses, and yaw symmetry is 0.92 to 0.96 against 0.68. % scorecard: enclosed_air_ratio, yaw_symmetry
The target is an exterior view: it fixes facades, roof and materials and says nothing about rooms, doors or stairs. % dream.yaml: dreams [iso]
The appearance gain and the structural loss are consistent with a builder that spends its edits on what the critic sees.
\TODO{test the interpretation with a cutaway or plan dream at n=256 (recipes dream-multiview, blueprint exist; T42 found the plan dream redraws the exterior).}

\subsection{Template collapse: every agentic arm fails the near-twin gate}

%TL;DR A fifth to a third of each agentic set has a near-copy inside the same set, against none of the real houses, so all three arms are disqualified.
Within each set, 31.6\% of \textsc{Classic-Matched} builds, 30.1\% of \textsc{Dream} builds and 21.1\% of \textsc{Dream-Caption} builds have a near-twin, against 0.0\% of real houses. % scorecard novelty: voxel_twin_rate
The gate disqualifies all three, at 15.8, 15.0 and 10.5 real spreads above the real rate. % scorecard leaderboard DQ reasons
Diversity agrees: 0.833, 0.786 and 0.790 against 0.902 for real houses. % scorecard voxel_diversity
On inspection, \textsc{Classic-Matched} draws most prompts as the same oak cottage with a gabled roof. % T51 (by eye)
The caption lowers the twin rate by 0.105 against the control, inside overlapping intervals ($[0.143, 0.286]$ against $[0.241, 0.391]$). % scorecard
Both trained models pass the near-twin gate (0.000 and 0.008). % scorecard

%TL;DR Before the twin gate existed, BlockScore ranked Dream-Caption above the other agentic arms and an LLM baseline, which shows why the gate is required.
Before the near-twin gate entered the protocol, BlockScore ranked \textsc{Dream-Caption} (155.8) above \textsc{Dream} (177.2) and \textsc{Classic-Matched} (181.9), and above an LLM block-serialization baseline (167.9). % T64 (n=128 card, 2026-10-23), before T91 (2026-11-04)
A leaderboard without the gate rewards an agent that repeats one house. % T91: ungated vanilla templates 202.95 / 194.98 ahead of every BlockGen model
No memorization gate fires on an agentic arm: DINO-NN percentiles are 0.460 to 0.591 against 0.536 for real houses. % scorecard dino_nn_percentile

\subsection{Trained models still lead}

%TL;DR The 20M trained model beats the best agentic arm on both co-primary metrics, and the 5M model beats Dream and Classic-Matched on geometry.
The 20M legacy voxel model beats \textsc{Dream-Caption} on Geom-KID by 6.65 ($[5.24, 8.08]$, $p_{\text{Holm}}{=}0.014$) and on MV-KID by 0.058 ($[0.030, 0.083]$, $p_{\text{Holm}}{=}0.014$). % scorecard h2h: recanon20m_s0 vs dream-caption -6.649, -0.058
The 5M \texttt{canon\_oriented} model beats \textsc{Dream} by 3.64 and \textsc{Classic-Matched} by 4.11 on Geom-KID (both $p_{\text{Holm}}{=}0.014$) and ties \textsc{Dream-Caption} on both metrics. % scorecard h2h: canon_oriented vs dream -3.638, vs classic -4.106; groups d/de, d/d
\texttt{canon\_oriented} is the only arm in \cref{tab:main} with a BlockScore (150.5). % scorecard leaderboard rank 2
On enterability \textsc{Classic-Matched} (0.677) exceeds both trained models (0.605 and 0.629). % scorecard traversal

%TL;DR The 20M model fails the memorization gate on this card and passes it on the n=256 card, which is the equal-n rule at work.
The 20M model fails the memorization gate here (DINO-NN percentile 0.359 against a real 0.536, 3.5 spreads) and passes every gate on an $n{=}256$ card. % scorecard leaderboard; T118 ("No gate fires on a model arm")
The gate's spread comes from a real sample of the card's $n$, so this verdict is one more reason to read the agentic and trained rows only through the pairwise tests. % docs/benchmark.md Protocols
The trained arms are scored at $n{=}256$ and the agentic arms at $n{=}133$; KID is unbiased in $n$, but the tests mix sample sizes. % scorecard n column
\TODO{n=256: rerun this comparison with all arms at n=256.}

\subsection{Cost}

%TL;DR A loop build costs 3.8 to 4.6 times a one-shot build, and the dream arms cost within 11% of the compute-matched control.
The three loop arms cost \$0.040 to \$0.048 per build with 3.5 to 4.6 model calls, against \$0.010 and one call for \textsc{One-Shot}, a factor of 3.8 to 4.6. % scorecard cost; run summaries n_llm_calls; oneshot paid $1.3892/133 = 0.01045; 0.040/0.01045=3.8, 0.048/0.01045=4.6
\textsc{Dream} costs 11\% less than \textsc{Classic-Matched} and \textsc{Dream-Caption} 7\% more. % 0.040/0.045, 6.385/5.968
Completion tokens per build are 14{,}600 for \textsc{Dream}, 17{,}962 for \textsc{Dream-Caption} and 20{,}315 for \textsc{Classic-Matched}, against 4{,}958 for \textsc{One-Shot}. % scorecard cost completion_tokens; oneshot summary.md
At these prices the full $n{=}133$ run of an arm costs \$5.30 to \$6.39. % run summaries total_would_cost_usd: 5.299, 5.968, 6.385
The trained models cost no API calls at sampling time. \TODO{add GPU-hours for training canon\_oriented and the 20M legacy model (T111, T118) so the cost comparison is complete.}

\subsection{Controls and the $n{=}256$ run}

\begin{table}[t]
\centering
\footnotesize
\setlength{\tabcolsep}{3pt}
\caption{\textbf{Controls and the $n{=}256$ card.} Each row isolates one component. All rows use the 256 prompts of \cref{sec:setup}.}
\label{tab:controls}
\resizebox{\columnwidth}{!}{%
\begin{tabular}{@{}llcc@{}}
\toprule
Arm & Isolates & MV-KID & Geom-KID \\
\midrule
\textsc{One-Shot} & the loop & \UNMEASURED{} & \UNMEASURED{} \\
\textsc{Classic-Matched} & (control) & \UNMEASURED{} & \UNMEASURED{} \\
\textsc{Dream} & the image & \UNMEASURED{} & \UNMEASURED{} \\
\textsc{Dream-Caption} & the caption & \UNMEASURED{} & \UNMEASURED{} \\
\textsc{Dream-Caption-Shuffled} & correct caption & \UNMEASURED{} & \UNMEASURED{} \\
\textsc{Text-Spec} & the picture & \UNMEASURED{} & \UNMEASURED{} \\
\textsc{Best-of-4} & more samples & \UNMEASURED{} & \UNMEASURED{} \\
\texttt{inspect-lessons} & text feedback & \UNMEASURED{} & \UNMEASURED{} \\
Gemma 4 12B (local) & the model & \UNMEASURED{} & \UNMEASURED{} \\
\bottomrule
\end{tabular}}
\end{table}
% State of outputs/logs/agentic256.log at 2026-11-09 20:58 UTC (in progress, not waited for):
%   One-Shot test (133) done: run_20261109_130512_bench_houses32_oneshot_test, paid $1.3892, not yet benched.
%   One-Shot val (123) done: run_20261109_145342_bench_houses32_oneshot_val, paid $1.3107, not yet benched.
%   Agentic val (123 x [dream-caption, dream, classic-matched]): run_20261109_165950_bench_houses32_agentic_val,
%   dream-caption at 86/123; dream and classic-matched val not started.
%   The test-split builds of the three loop arms already exist (T51 runs, dream-caption133).

%TL;DR The control table fixes in advance which comparison answers which question, so the n=256 card reads off directly.
\Cref{tab:controls} fixes the comparisons the $n{=}256$ card answers. \TODO{n=256}
\textsc{Dream-Caption} against \textsc{Dream-Caption-Shuffled} measures what a correct description adds beyond more text; if they tie, the gain is tokens. % docs/agentic-suite.md 5.6
\textsc{Dream-Caption} against \textsc{Text-Spec} measures what the picture adds beyond a description.
\textsc{Dream} against \textsc{Best-of-4} at matched cost measures whether the image beats sampling four times. % best-of-4.yaml comment; docs 9
\texttt{inspect-lessons} against \textsc{Classic-Matched} measures whether text geometry feedback fixes the enterability gap of \cref{tab:inspect}, which a render critic does not see. % inspection.py calibrated on 400 real houses (mean 0.93), memory note agentic-suite
The local Gemma arm measures whether the findings depend on the closed model. % T47, T48 pilots n=6

\subsection{Limitations}

%TL;DR One run per prompt, one builder, one image model, one building type and one canvas size per arm bound every claim above.
Every arm runs once per prompt with one builder (\texttt{gpt-5-mini}) and one image model (\texttt{gpt-image-1-mini}), so the findings carry no seed variance and no model variance. % summaries: n_builds = prompts; dream.yaml
The prompts are houses at $32^3$, and the $n{=}133$ card sits just above the protocol's minimum of 128 per arm. % docs/benchmark.md min_n 128
\textsc{One-Shot} ran on a $48^3$ canvas, unlike the loop arms. % oneshot config.json
The selectors differ between arms (CLIP text for the control, SigLIP image similarity for the dream arms), so a dream arm changes both its input and its selection signal. % recipes score fields
The titles name real GrabCraft builds that the builder may have seen during pretraining; the memorization gate checks proximity to our training split, not to the web. % scorecard dino_nn_percentile is against the training split
The $n{=}256$ prompt set adds validation titles while the reference stays the test split. % prompts files; scorecard reference 399

\section{Conclusion}
\label{sec:conclusion}

%TL;DR Visual imagination buys a program-writing agent appearance, through the caption of the image, and buys no structure.
We ask what a target image buys an LLM that writes build programs, and we answer with a compute-matched control on 133 held-out house prompts. % T51 / scorecard n=133
The image alone does not separate from the control on either co-primary metric. % scorecard groups
A caption of the image does: \textsc{Dream-Caption} improves MV-KID by 0.186 over the control and is not separated from a 5M trained model on appearance. % scorecard h2h; canon_oriented group d
Geometry stays tied across all three loop arms, and structure gets worse: a player enters half of the dream houses against two thirds of the control houses and 93\% of real ones. % scorecard traversal
Every agentic arm repeats itself often enough to fail the near-twin gate, and a 20M trained voxel model still leads on both metrics. % scorecard leaderboard, h2h
Imagination, in this setting, is a lever on appearance and not on structure.

%TL;DR The evaluation protocol is what makes the negative half of the result visible: without Geom-KID, traversal and the twin gate, the dream loop looks like progress.
Each negative finding comes from a part of the protocol that a render metric lacks. % anonymous2027blindness
Geom-KID shows the geometry tie, the traversal probe and the inspector show the entry failures, and the near-twin gate shows the template collapse that an ungated BlockScore rewarded. % T63, T64, T91, T95
An agentic method that reports only render similarity to its own target reports the one axis on which the image helps.

%TL;DR Open questions: interior targets, symbolic feedback, sampling controls, local models, and more seeds.
Five questions remain open. % ideas, research tracks
First, whether a target that shows the interior, such as a cutaway or a floor plan, moves structure the way the exterior moves appearance; in our pilots, image models redraw the exterior when asked for a plan. % T42
Second, whether symbolic geometry feedback (\texttt{inspect-lessons}) closes the enterability gap that a render critic does not see. % inspect-lessons.yaml, T95
Third, whether the caption's gain survives the shuffled-caption, text-spec and best-of-4 controls of \cref{tab:controls}. \TODO{n=256}
Fourth, whether the findings hold for an open local model with an open image generator. % T47, T48 pilots
Fifth, how large the seed and model variance is, which a single run per prompt does not measure.


{
    \small
    \bibliographystyle{ieeenat_fullname}
    \bibliography{references,references_new}
}

\end{document}
